%% parts/p01/formal-work-conventions.tex --- compressed usable core of the
%% verified “Levels” and “Writing” chapters.  It is lookup material; exact
%% grammars, semantic recursions, and rule displays remain in their home parts.
\chapter{Levels, Consequence, and Formal-Work Conventions}
\chaptermark{Formal-Work Conventions}
\label{sec:p01-formal-work}

These clauses fix distinctions used throughout the handbook.  They are
schematic until a system chapter declares its exact grammar, calculus, and
semantic data; no entry below substitutes for those system-specific
definitions.

\section{Levels and system components}
\label{sec:p01-levels-systems}

\begingroup\small
\begin{longtable}{@{}p{0.20\textwidth}p{0.51\textwidth}p{0.21\textwidth}@{}}
\caption{Language, system, and level vocabulary.}
\label{tab:p01-level-vocabulary}\\
\toprule
Term & Exact clause & Standing \\
\midrule
\endfirsthead
\toprule
Term & Exact clause & Standing \\
\midrule
\endhead
Formal language\index{formal language} & A set of finite expressions
determined from a specified alphabet by a specified grammar; the grammar
fixes each grammatical category. & abstract syntactic object \\
Calculus\index{calculus} & Specified initial derivations and inference rules
that inductively determine a derivability relation. & syntactic component \\
Semantics\index{semantics} & The complete permitted semantic data for a
language together with a satisfaction relation between those data and
formulas. & semantic component \\
Logical system\index{logical system} & An ordered triple
\(\langle\mathcal L,\mathcal C,\mathcal S\rangle\) of a language, a
calculus for it, and a semantics for it; the term is used only when all three
components are identified. & triple of components \\
Object language\index{object language} & The formal language under study; its
formulas are mathematical objects, and writing one does not assert it. &
object of study \\
Metalanguage\index{metalanguage} & The language in which the object language,
its calculi and semantics, and their metatheory are defined; it includes the
ambient mathematics and this prose. & language of analysis \\
Syntax\index{syntax} & The alphabet and grammar of a formal language together
with its operations on expressions defined without reference to
interpretation. & syntactic data of a language \\
Proof theory\index{proof theory} & Study of derivations, proof
transformations, and properties of calculi. & metalinguistic subject \\
\bottomrule
\end{longtable}
\endgroup

Changing natural deduction to a sequent calculus changes \(\mathcal C\)
without necessarily changing \(\mathcal L\) or \(\mathcal S\); changing
two-valued valuations to a three-valued matrix changes \(\mathcal S\).
Accordingly, \(\derives\) and \(\entails\) never name universal relations:
their registered subscripts are mandatory when more than one calculus or
semantics is active.  The object/metalanguage distinction is functional, not
typographical, and becomes substantive when syntax is internally coded
\cite{Tarski1983}.

\begin{hbwarning}{def:p01-level-confusion}{Use, mention, and logical level}
Normally \(\varphi,\psi\) are metavariables ranging over formulas, not
symbols of the object language.  If \(\to\) belongs to the grammar,
\(\varphi\to\psi\) denotes the resulting family of formulas.  The lifted
claims \(\{\varphi\}\entails_{\mathcal S}\psi\) and
\(\{\varphi\}\derives_{\mathcal C}\psi\) are metalinguistic and require
the intended universal quantification over formulas.  Never write
\(\varphi\entails\psi\) as though a formula metavariable were already a
singleton premise set.  Quotation or a code mentions an expression; an
unquoted inscription inside the active grammar uses it.
\end{hbwarning}

\section{Consequence, satisfaction, and arguments}
\label{sec:p01-consequence-vocabulary}

\begingroup\small
\begin{longtable}{@{}p{0.22\textwidth}p{0.53\textwidth}p{0.17\textwidth}@{}}
\caption{Semantic and proof-theoretic vocabulary.}
\label{tab:p01-consequence-vocabulary}\\
\toprule
Term & Exact clause & Kind \\
\midrule
\endfirsthead
\toprule
Term & Exact clause & Kind \\
\midrule
\endhead
Satisfaction\index{satisfaction} & The relation between a formula and every
piece of semantic data required by the active semantics; e.g.
\(\mathcal M,s\entails\varphi\) in first-order logic. & semantic relation \\
Truth in a structure\index{truth!in a structure} & For a sentence
\(\sigma\), truth in \(\mathcal M\) means
\(\mathcal M,s\entails\sigma\) for every variable assignment \(s\).
The later assignment-independence theorem licenses the abbreviation
\(\mathcal M\entails\sigma\).  For a formula with free occurrences, the
assignment may not be suppressed. & semantic property \\
Model\index{model} & \(\mathcal M\) is a model of \(\sigma\) iff
\(\mathcal M\entails\sigma\), and a model of \(T\) iff it satisfies every
member of \(T\).  A bare structure is not a model without a stated sentence
or theory. & relation of a structure to content \\
Semantic consequence\index{semantic consequence} &
\(\Gamma\entails_{\mathcal S}\varphi\) iff every permitted semantic datum
satisfying every member of \(\Gamma\) satisfies \(\varphi\). & semantic
relation \\
Derivability\index{derivability} &
\(\Gamma\derives_{\mathcal C}\varphi\) iff a \(\mathcal C\)-derivation of
\(\varphi\) exists with no undischarged assumptions outside \(\Gamma\). &
syntactic relation \\
Validity\index{validity} & \(\varnothing\entails_{\mathcal S}\varphi\). &
semantic property \\
Theorem of a calculus\index{theorem!of a calculus} &
\(\varnothing\derives_{\mathcal C}\varphi\). & syntactic property \\
Soundness\index{soundness} & For all \(\Gamma,\varphi\),
\(\Gamma\derives_{\mathcal C}\varphi\Rightarrow
\Gamma\entails_{\mathcal S}\varphi\). & derivability to consequence \\
Completeness\index{completeness!of a calculus} & For all
\(\Gamma,\varphi\),
\(\Gamma\entails_{\mathcal S}\varphi\Rightarrow
\Gamma\derives_{\mathcal C}\varphi\). & consequence to derivability \\
Argument\index{argument} & An ordered pair \(\langle\Gamma,\varphi\rangle\)
of premises and conclusion in one language; semantically valid iff
\(\Gamma\entails_{\mathcal S}\varphi\) for the named semantics. & syntactic
data with a semantic property \\
Countermodel\index{countermodel} & When every member of
\(\Gamma\cup\{\varphi\}\) is a sentence, a countermodel is a structure
\(\mathcal M\) satisfying every member of \(\Gamma\) and not
\(\varphi\).  For formulas with free occurrences, name a structure and
assignment instead. & sentence-only witness against consequence \\
\bottomrule
\end{longtable}
\endgroup

\begin{hbwarning}{def:p01-direction-check}{Direction check}
Soundness rules out derivations of semantically invalid conclusions;
completeness supplies derivations for semantic consequences.  Neither
direction is about an unnamed calculus or semantic class.
\end{hbwarning}

\begin{hbmethod}{def:p01-formalizing-argument}{Formalizing an argument}
\index{formalization!argument}
Work in this order: state the intended domain and the conditions for counting
an expression as a premise or conclusion; declare every nonlogical symbol,
with its arity or sort and intended interpretation; resolve pronouns, tense,
modality, plurality, restricted quantification, uniqueness, and ambiguity;
translate each premise and the conclusion separately while preserving
quantifier scope; state the treatment of free variables; and test the
translation against intended positive and negative cases.  Only then fix the
semantic class or calculus and ask for consequence or derivability.  For
sentence arguments, search for a countermodel when consequence is doubtful;
for formulas with free occurrences, search for a structure and assignment.
In particular, \(\forall x\exists y\,Rxy\) and
\(\exists y\forall x\,Rxy\) impose different model conditions.
\end{hbmethod}

\section{Definition and theorem audit}
\label{sec:p01-definition-proof-audit}

A dependency-correct development introduces, in order: ambient sets and
metalinguistic conventions; signature and arities; term and formula grammar;
free variables and substitution; structures and assignments; evaluation and
satisfaction; calculus and derivations; their properties; and finally the
connections between calculus and semantics.  A stated pedagogical reordering
does not license definition before use.

\begin{hbmethod}{def:p01-definition-audit}{Definition audit}
\index{definition!audit}
Record the item's type or standing; its previously defined range; all
quantifiers; function domains and codomains; arities and binding positions;
totality or domain of definition; a recursion measure or grammar; and, for a
quotient representative, independence of representative.  A new technical
name additionally requires necessity, an exact eliminable definition or
explicit axiom, standing, nonredundancy, and an actual later use.  Until then,
describe the owed artifact without naming a hypothetical answer.
\end{hbmethod}

A theorem statement identifies every quantified object and hypothesis,
including finiteness, countability, consistency, effectiveness, and choice;
the exact language, calculus, and semantic class; the direction of each
implication; and whether its objects are formulas, sentences, theories,
structures, or derivations.  ``Clearly,'' ``obviously,'' and ``routine'' do
not discharge a proof obligation.

\begin{hbmethod}{def:p01-proof-plan}{Proof plan before proof}
\index{proof plan}
Record the fully quantified target, governing definition or theorem, proof
method, induction parameter or minimal counterexample, existential witness,
and maintained invariant.  Remove the plan only if the proof remains
auditable.
\end{hbmethod}

\begingroup\small
\begin{longtable}{@{}p{0.19\textwidth}p{0.71\textwidth}@{}}
\caption{Proof-method obligations.}
\label{tab:p01-proof-methods}\\
\toprule
Method & Required audit \\
\midrule
\endfirsthead
\toprule
Method & Required audit \\
\midrule
\endhead
Direct\index{proof!direct} & Assume the antecedent and derive the consequent
at the same level; a semantic-to-syntactic change requires a prior theorem. \\
Contrapositive\index{proof!contrapositive} & Legitimate in classical
metalogic; the corresponding object-language principle is not generally
intuitionistically derivable. \\
Contradiction\index{proof!by contradiction} & Classical derivation of
\(\bot\) from \(\neg A\) proves \(A\); intuitionistically it proves only
\(\neg\neg A\) without a further principle. \\
Induction\index{proof!by induction} & State domain, predicate, bases, and
step.  Structural induction follows grammar; derivation induction follows the
final rule; complexity induction follows its defined measure. \\
Construction\index{proof!by construction} & Define the witness and verify
every property; a model includes its domain, all interpretations, and a
satisfaction proof. \\
Counterexample\index{counterexample} & Exhibit an object in the quantified
domain and verify failure; with free variables, semantic refutation supplies
both a structure and an assignment. \\
\bottomrule
\end{longtable}
\endgroup

\section{Prose, notation, and rule references}
\label{sec:p01-formula-prose}

Displayed formulas are punctuated parts of sentences.  Introduce each symbol
before use and justify every nontrivial line of a transformation.  Under the
target glyph rulings,
\[
 \neg(P\to Q)\Leftrightarrow\neg(\neg P\lor Q)
 \Leftrightarrow P\land\neg Q.
\]
Here \(\Leftrightarrow\) is formula equivalence,
\(\leftrightarrow\) is the object-language biconditional, and \(\equiv\)
is reserved for elementary equivalence.  Rule displays occur only in
\texttt{registry/calculi/}: body prose cites the relevant
\texttt{sheet:...} or \texttt{rule:...} label, whose registry row carries
the name and every side condition.  On mathematical prose and notation see
\cite{Halmos1970,KnuthLarrabeeRoberts1989,AMSAuthorHandbook}.

\section{Error taxonomy}
\label{sec:p01-error-taxonomy}
\index{error taxonomy}

\begingroup\small
\begin{longtable}{@{}p{0.20\textwidth}p{0.34\textwidth}p{0.37\textwidth}@{}}
\caption{Formal-work error taxonomy and repairs.}
\label{tab:p01-error-taxonomy}\\
\toprule
Error & Diagnostic question & Repair \\
\midrule
\endfirsthead
\toprule
Error & Diagnostic question & Repair \\
\midrule
\endhead
Level & Is a formula confused with a metalinguistic statement? & Mention it
by quotation or metavariable and state the metalinguistic relation. \\
Grammar & Is the expression generated by the active grammar? & Parse it from
that grammar or correct it. \\
Binding & Which binder governs each occurrence? & Rename bound variables and
restore parentheses. \\
Capture & Does substitution bind a formerly free occurrence? & Rename the
conflicting bound variable first. \\
Semantic & Is truth asserted without all required semantic data? & Supply
every semantic parameter. \\
Rule & Is the rule in the active calculus, with all side conditions met? &
Cite its registry row and verify each condition. \\
Scope & Does an operator govern the intended subformula? & Add parentheses or
rewrite with explicit restricted quantification. \\
Metatheorem & Is soundness, completeness, consistency, or decidability used in
the wrong direction? & Write the quantified implication and match every
hypothesis. \\
Strength & Does the argument establish less than the result states? &
Strengthen the proof or weaken and relabel the statement; never upgrade its
status. \\
\bottomrule
\end{longtable}
\endgroup

\section{Further reading}
\label{sec:p01-further-reading}

For the mathematical substrate see \cite{Jech2003}; for the
syntax--semantics--calculus distinctions see
\cite{Enderton2001,EbbinghausFlumThomas1994,OpenLogic2026}; for levels and
truth see \cite{Tarski1983}; and for exposition see
\cite{Halmos1970,KnuthLarrabeeRoberts1989,AMSAuthorHandbook}.
